% ============================================================================
% DOCUMENTO: Guia_Evaluacion_ATRIL.tex
% PROYECTO:  Plataforma ATRIL v5.0 (MVP) — FundaOboeHJK & USB
% MATERIA:   PS-1115: Sistemas de Información I (Trimestre Sep-Dic 2026)
% LAB 6:     Pruebas de Software — Guía de Pruebas y Coevaluación (Equipo 1)
% COMPILAR:  Overleaf (pdfLaTeX / XeLaTeX) — Exactamente 2 páginas
% ============================================================================
\documentclass[10pt,letterpaper]{article}

% --- Paquetes Obligatorios y de Maquetación ---
\usepackage[utf8]{inputenc}
\usepackage[T1]{fontenc}
\usepackage[spanish,es-tabla]{babel}
\usepackage[scaled=0.94]{helvet}
\renewcommand{\familydefault}{\sfdefault}

\usepackage[
  letterpaper,
  top=1.25cm,
  bottom=1.65cm,
  left=1.40cm,
  right=1.40cm,
  headheight=14pt,
  footskip=20pt
]{geometry}

\usepackage[table]{xcolor}
\usepackage{tabularx}
\usepackage{array}
\usepackage[most]{tcolorbox}
\usepackage{listings}
\usepackage{enumitem}
\usepackage{fancyhdr}
\usepackage{amssymb}
\usepackage[hidelinks]{hyperref}

% --- Paleta Institucional ATRIL v5.0 (src/constants/theme.ts) ---
\definecolor{atrilPrimary}{HTML}{0047AB}      % Azul Rey (#0047AB)
\definecolor{atrilPrimaryHover}{HTML}{003580} % Azul Rey Oscuro (#003580)
\definecolor{atrilNavy}{HTML}{0A192F}         % Azul Marino Profundo (#0A192F)
\definecolor{atrilBg}{HTML}{F4F7FC}           % Fondo Suave (#F4F7FC)
\definecolor{atrilSurface}{HTML}{FFFFFF}      % Blanco (#FFFFFF)
\definecolor{atrilSurfaceSub}{HTML}{EBF2FA}   % Superficie Secundaria (#EBF2FA)
\definecolor{atrilGold}{HTML}{FFBF00}         % Acento Dorado (#FFBF00)

% --- Estados Semánticos (WCAG AA) ---
\definecolor{atrilSuccess}{HTML}{047857}
\definecolor{atrilSuccessTint}{HTML}{D1FAE5}
\definecolor{atrilSuccessText}{HTML}{065F46}
\definecolor{atrilWarnTint}{HTML}{FEF3C7}
\definecolor{atrilWarnText}{HTML}{92400E}
\definecolor{atrilDanger}{HTML}{C81E1E}
\definecolor{atrilDangerTint}{HTML}{FEE2E2}
\definecolor{atrilDangerText}{HTML}{991B1B}
\definecolor{atrilTextBody}{HTML}{1E293B}
\definecolor{atrilTextMuted}{HTML}{334155}
\definecolor{atrilBorder}{HTML}{CBD5E1}
\definecolor{termComment}{HTML}{94A3B8}

% --- Configuración Global de Párrafo y Tablas ---
\setlength{\parindent}{0pt}
\setlength{\parskip}{0pt}
\renewcommand{\arraystretch}{1.32}
\newcolumntype{Y}{>{\raggedright\arraybackslash}X}

% --- Configuración de Encabezado y Pie de Página (fancyhdr) ---
\pagestyle{fancy}
\fancyhf{}
\renewcommand{\headrulewidth}{0pt}
\renewcommand{\footrulewidth}{0.6pt}
\renewcommand{\footrule}{\hbox to\headwidth{\color{atrilBorder}\leaders\hrule height \footrulewidth\hfill}}

\fancyfoot[L]{%
  \scriptsize\color{atrilTextMuted}%
  \textbf{Equipo ATRIL (FundaOboeHJK)} --- José Daniel Torbet, Samuel Sanabria, Nadheska García, Ernesto Barreto%
}
\fancyfoot[R]{%
  \scriptsize\color{atrilTextMuted}%
  \ifnum\value{page}=1%
    \textbf{Página 1 de 2. Guía de Evaluación}%
  \else%
    \textbf{Página 2 de 2. Documento Oficial de Coevaluación}%
  \fi%
}

% --- Macros Visuales Corporativos ---
\newcommand{\secbadge}[2]{%
  \vspace{0.22cm}%
  \noindent\colorbox{atrilNavy}{\strut\hspace{4pt}\textcolor{atrilGold}{\textbf{\small #1}}\hspace{4pt}}%
  \hspace{6pt}{\large\bfseries\color{atrilNavy}#2}\par\vspace{0.12cm}%
}

\newcommand{\pillRolAdmin}{\colorbox{atrilWarnTint}{\textcolor{atrilWarnText}{\textbf{\scriptsize ADMIN}}}}
\newcommand{\pillRolDocente}{\colorbox{atrilSurfaceSub}{\textcolor{atrilPrimary}{\textbf{\scriptsize DOCENTE}}}}
\newcommand{\pillRolEstudiante}{\colorbox{atrilSurfaceSub}{\textcolor{atrilNavy}{\textbf{\scriptsize ESTUDIANTE}}}}

\newcommand{\pillVerificado}{\colorbox{atrilSuccessTint}{\textcolor{atrilSuccessText}{\textbf{\scriptsize VERIFICADO}}}}
\newcommand{\pillPendiente}{\colorbox{atrilWarnTint}{\textcolor{atrilWarnText}{\textbf{\scriptsize PENDIENTE}}}}

\newcommand{\codeinline}[1]{\colorbox{atrilSurfaceSub}{\texttt{\footnotesize\textcolor{atrilNavy}{#1}}}}

% --- Configuración de Listings para Bloques de Terminal ---
\lstset{
  basicstyle=\ttfamily\footnotesize\color{white},
  breaklines=true,
  columns=fullflexible,
  keepspaces=true,
  showstringspaces=false,
  commentstyle=\color{termComment}\itshape,
  literate=
    {á}{{\'a}}1 {é}{{\'e}}1 {í}{{\'i}}1 {ó}{{\'o}}1 {ú}{{\'u}}1
    {Á}{{\'A}}1 {É}{{\'E}}1 {Í}{{\'I}}1 {Ó}{{\'O}}1 {Ú}{{\'U}}1
    {ñ}{{\~n}}1 {Ñ}{{\~N}}1
}

\begin{document}
\color{atrilTextBody}

% ============================================================================
% PÁGINA 1: GUÍA DE PRUEBAS Y COEVALUACIÓN
% ============================================================================

% --- Encabezado Página 1 ---
\noindent
\begin{tabularx}{\textwidth}{@{}Y r@{}}
  \begin{minipage}[t]{\linewidth}
    {\LARGE\bfseries\color{atrilNavy} GUÍA DE PRUEBAS Y COEVALUACIÓN}\\[3pt]
    {\normalsize\bfseries\color{atrilPrimary} Plataforma ATRIL --- Red de Conexión Docente y Chat Interno (FundaOboeHJK)}
  \end{minipage}
  &
  \begin{minipage}[t]{0.36\textwidth}\raggedleft\footnotesize\color{atrilTextMuted}
    \textbf{\color{atrilNavy}Universidad Simón Bolívar}\\
    PS-1115: Sistemas de Información I\\
    Laboratorio 6: Pruebas de Software\\
    Trimestre Sep--Dic 2026
  \end{minipage}
\end{tabularx}

\vspace{0.14cm}
{\color{atrilPrimary}\hrule height 2pt}
{\color{atrilGold}\hrule height 1pt}
\vspace{0.18cm}

% --- Bloque de Acceso y Despliegue ---
\begin{tcolorbox}[
  colback=atrilBg,
  colframe=atrilPrimary,
  boxrule=0.9pt,
  arc=4pt,
  left=8pt,
  right=8pt,
  top=5pt,
  bottom=5pt
]
  \textbf{\small\color{atrilNavy} Arquitectura Operativa ATRIL --- Entorno Multiplataforma (Expo Web / Móvil + Supabase BaaS)}
  {\footnotesize\color{atrilTextBody}
  La versión actual implementa el flujo institucional tutelado por \textbf{FundaOboeHJK}: vitrina pública protegida por candado, postulación docente y verificación estudiantil (adultos y menores con representante legal), consola Superadmin y mensajería interna en tiempo real.}
\end{tcolorbox}

\vspace{0.10cm}

% Tarjetas 2x2 de Acceso y Despliegue
\noindent
\begin{minipage}[t]{0.49\textwidth}
  \begin{tcolorbox}[colback=white,colframe=atrilBorder,boxrule=0.7pt,arc=3pt,left=7pt,right=7pt,top=4pt,bottom=4pt,height=1.30cm,valign=center]
    {\scriptsize\bfseries\color{atrilTextMuted} APLICACIÓN EN VIVO / EXPO WEB}\\[1pt]
    {\scriptsize\bfseries\color{atrilPrimary}\url{https://github.com/Nadhe15/Atril-/tree/feature/auth-y-reserva-modular}}\\[1pt]
    {\scriptsize\color{atrilTextMuted} URL Web: \texttt{https://visionary-cocada-c37d8f.netlify.app} \textbar\ Local: \texttt{npx expo start --web}}
  \end{tcolorbox}
\end{minipage}\hfill
\begin{minipage}[t]{0.49\textwidth}
  \begin{tcolorbox}[colback=white,colframe=atrilBorder,boxrule=0.7pt,arc=3pt,left=7pt,right=7pt,top=4pt,bottom=4pt,height=1.30cm,valign=center]
    {\scriptsize\bfseries\color{atrilTextMuted} BACKEND / BaaS (AUTENTICACIÓN Y TIEMPO REAL)}\\[1pt]
    {\footnotesize\bfseries\color{atrilPrimary} Supabase Auth + PostgreSQL + Realtime (RLS)}\\[1pt]
    {\scriptsize\color{atrilTextMuted} Con caché local (\texttt{adminService} / \texttt{chatService})}
  \end{tcolorbox}
\end{minipage}

\vspace{0.08cm}

\noindent
\begin{minipage}[t]{0.49\textwidth}
  \begin{tcolorbox}[colback=white,colframe=atrilBorder,boxrule=0.7pt,arc=3pt,left=7pt,right=7pt,top=4pt,bottom=4pt,height=1.12cm,valign=center]
    {\scriptsize\bfseries\color{atrilTextMuted} REPOSITORIO DE CÓDIGO FUENTE (GITHUB)}\\[1pt]
    {\scriptsize\bfseries\color{atrilPrimary}\url{https://github.com/Nadhe15/Atril-/tree/feature/auth-y-reserva-modular}}
  \end{tcolorbox}
\end{minipage}\hfill
\begin{minipage}[t]{0.49\textwidth}
  \begin{tcolorbox}[colback=white,colframe=atrilBorder,boxrule=0.7pt,arc=3pt,left=7pt,right=7pt,top=4pt,bottom=4pt,height=1.12cm,valign=center]
    {\scriptsize\bfseries\color{atrilTextMuted} CONTRASEÑA UNIVERSAL / ACCESO DEMO}\\[1pt]
    {\footnotesize\bfseries\color{atrilNavy} Botonera Demo 1-Clic (\texttt{AuthLandingScreen}) \textbar\ Mín. 6 caract.}
  \end{tcolorbox}
\end{minipage}

% --- Sección 1: Cuentas de Acceso Preconfiguradas ---
\secbadge{1}{Cuentas de Acceso Preconfiguradas (De acceso rápido para la Demo)}
{\footnotesize\color{atrilTextMuted}
Cuentas reales preconfiguradas en \texttt{src/screens/AuthLandingScreen.tsx} y semillas en \texttt{src/services/adminService.ts} / Supabase:}

\vspace{0.10cm}
\noindent
{\footnotesize
\begin{tabularx}{\textwidth}{|>{\ttfamily\scriptsize}l|Y|c|c|l|}
  \hline
  \rowcolor{atrilNavy}
  \textcolor{white}{\textbf{\textsf{Usuario / Correo}}} &
  \textcolor{white}{\textbf{\textsf{Nombre Completo (Cátedra / Condición)}}} &
  \textcolor{white}{\textbf{\textsf{Rol Sistema}}} &
  \textcolor{white}{\textbf{\textsf{Estado Verif.}}} &
  \textcolor{white}{\textbf{\textsf{Acceso / Contraseña}}} \\
  \hline
  \rowcolor{atrilBg}
  virginia.admin@fundaoboe.org.ve &
  \textbf{Admin} (Superadmin Fundaoboe) &
  \pillRolAdmin &
  \pillVerificado &
  \codeinline{|Demo| Admin} \\
  \hline
  \rowcolor{white}
  jg.sanabria@fundaoboe.org &
  \textbf{Mtro. José Gregorio Sanabria} (\emph{Oboe Sinfónico}) &
  \pillRolDocente &
  \pillVerificado &
  \codeinline{Demo: Profesor Verificado} \\
  \hline
  \rowcolor{atrilBg}
  samuel.sanabria@usb.ve &
  \textbf{Samuel Sanabria} (\emph{Estudiante Adulto \textbullet\ USB}) &
  \pillRolEstudiante &
  \pillVerificado &
  \codeinline{Demo: Estudiante Verificado} \\
  \hline
  \rowcolor{white}
  carlos.reveron@conservatorio.ve &
  \textbf{Prof. Carlos Reverón} (\emph{Guitarra Clásica}) &
  \pillRolDocente &
  \pillPendiente &
  \codeinline{Demo: Usuario Pendiente} \\
  \hline
  \rowcolor{atrilBg}
  carolina.mendoza@correo.com &
  \textbf{Mateo Mendoza} (\emph{Menor \textbullet\ Rep: Carolina Mendoza}) &
  \pillRolEstudiante &
  \pillPendiente &
  \codeinline{Semilla Admin (usr-pend-02)} \\
  \hline
\end{tabularx}
}

\vspace{0.08cm}
{\scriptsize\color{atrilTextMuted}
\textbf{Catálogo Institucional Precargado (\texttt{teacherCatalogService.ts}):} 5 docentes verificados: \emph{Mtro. José Gregorio Sanabria} (Oboe Sinfónico), \emph{Prof. Valentina Meneses} (Fagot), \emph{Mtro. Leonardo Montilla} (Guitarra Clásica y Cuatro), \emph{Prof. Mariana Salcedo} (Violín) y \emph{Prof. Adriana Villamizar} (Canto).}

% --- Sección 2: Procedimiento de Verificación en la Interfaz ---
\secbadge{2}{Procedimiento de Verificación en la Interfaz}

\begin{itemize}[leftmargin=0.55cm, itemsep=3.5pt, topsep=2pt, parsep=0pt]
  \item[\textcolor{atrilPrimary}{\textbf{$\checkmark$}}] \footnotesize
  \textbf{1. Pantalla de Bienvenida (\texttt{SplashScreen}) y Vista Externa Pública (\texttt{AuthLandingScreen}):}
  Al abrir la app se despliega la portada institucional de la \emph{Fundación Venezolana}. Pulse \codeinline{Ingresar a la Plataforma $\rightarrow$} y explore las pestañas superiores \textbf{Quiénes Somos} y \textbf{Profesores (Vista Externa)}, comprobando que los perfiles públicos muestran su cátedra pero mantienen los datos de contacto bloqueados por candado (\codeinline{Inicia sesión para ver perfil}).

  \item[\textcolor{atrilPrimary}{\textbf{$\checkmark$}}] \footnotesize
  \textbf{2. Registro y Estado Pendiente (\texttt{PendingVerificationScreen}):}
  En la pestaña de autenticación, seleccione \codeinline{Postularse / Verificarse}. Pruebe \textbf{Postularse como Profesor} (eligiendo una de las 8 cátedras oficiales) o \textbf{Verificarse como Estudiante} (probando la opción \emph{Menor de 18 Con Representante}, que valida nombre y cédula del tutor). Verifique que tras registrarse la cuenta queda bloqueada en \texttt{PendingVerificationScreen} en espera de acreditación externa por FundaOboe (\texttt{acreditacion@fundaoboe.org.ve}).

  \item[\textcolor{atrilPrimary}{\textbf{$\checkmark$}}] \footnotesize
  \textbf{3. Accesos Rápidos de Prueba (Demo Toolbar):}
  En el modo \emph{Iniciar Sesión} de \texttt{AuthLandingScreen}, utilice la botonera inferior \textbf{Accesos Rápidos de Prueba (Demo)} para ingresar en 1 clic como \codeinline{Admin}, \codeinline{Profesor Verificado}, \codeinline{Estudiante Verificado} o \codeinline{Usuario Pendiente}. Dentro de cualquier sesión activa, pulse \codeinline{$\rightleftarrows$ Cambiar cuenta / Salir} en la barra superior para alternar de rol sin cerrar ni crear cuentas manualmente.

  \item[\textcolor{atrilPrimary}{\textbf{$\checkmark$}}] \footnotesize
  \textbf{4. Panel Superadmin (\texttt{AdminDashboardScreen}):}
  Ingrese como \codeinline{Admin }, filtre los usuarios registrados por estado (\emph{Todos, Pendientes, Verificados, Suspendidos}) y por rol (\emph{Profesor, Estudiante}), y pruebe los botones de acción directa \codeinline{$\checkmark$ Aprobar / Verificar} y \codeinline{$\times$ Suspender / Desactivar}. Al aprobar a \textbf{Prof. Carlos Reverón} y volver a ingresar con el botón demo \codeinline{Usuario Pendiente}, compruebe que el bloqueo desaparece y accede como Profesor Verificado.

  \item[\textcolor{atrilPrimary}{\textbf{$\checkmark$}}] \footnotesize
  \textbf{5. Directorio de Profesores y Chat Interno (\texttt{StudentTeachersScreen}, \texttt{TeacherDetailScreen}, \texttt{ChatInboxScreen}, \texttt{ChatRoomScreen}):}
  Ingrese como \codeinline{Estudiante Verificado} (Samuel Sanabria), filtre profesores por área/cátedra (\emph{Oboe, Guitarra, Cuatro, Fagot, Violín, Canto, Armonía}), abra el perfil del \textbf{Mtro. José Gregorio Sanabria}, elija el área a estudiar y pulse \codeinline{Iniciar Chat con el Profesor $\rightarrow$}. Envíe un mensaje en \texttt{ChatRoomScreen} y alterne a la cuenta \codeinline{Profesor Verificado} (\texttt{ChatInboxScreen}) para verificar el intercambio bidireccional de mensajes.
\end{itemize}

\newpage

% ============================================================================
% PÁGINA 2: SUITE TÉCNICA Y PROTOCOLO DE DEFECTOS
% ============================================================================

% --- Encabezado Página 2 ---
\noindent
\begin{tabularx}{\textwidth}{@{}Y r@{}}
  \begin{minipage}[t]{\linewidth}
    {\LARGE\bfseries\color{atrilNavy} SUITE TÉCNICA Y PROTOCOLO DE DEFECTOS}\\[3pt]
    {\normalsize\bfseries\color{atrilPrimary} Pruebas Automatizadas, Verificación de Tipos e Inyección de Defectos (Rúbrica Lab 6)}
  \end{minipage}
  &
  \begin{minipage}[t]{0.36\textwidth}\raggedleft\footnotesize\color{atrilTextMuted}
    \textbf{\color{atrilNavy}Universidad Simón Bolívar}\\
    PS-1115: Sistemas de Información I\\
    Laboratorio 6: Pruebas de Software\\
    Trimestre Sep--Dic 2026
  \end{minipage}
\end{tabularx}

\vspace{0.14cm}
{\color{atrilPrimary}\hrule height 2pt}
{\color{atrilGold}\hrule height 1pt}
\vspace{0.18cm}

% --- Sección 3: Ejecución Local del Proyecto y Suite de Pruebas ---
\secbadge{3}{Ejecución Local del Proyecto y Suite de Pruebas}
{\footnotesize\color{atrilTextBody}
Comandos reales verificados sobre el repositorio \texttt{atril-app} (Expo SDK 58 + React Native 0.88 + TypeScript 6.0):}

\vspace{0.10cm}
\begin{tcolorbox}[
  colback=atrilNavy,
  colframe=atrilPrimary,
  boxrule=1pt,
  arc=4pt,
  left=8pt,
  right=8pt,
  top=5pt,
  bottom=5pt
]
\begin{lstlisting}
# 1. Clonar el repositorio en la rama feature/auth-y-reserva-modular y acceder al directorio
git clone -b feature/auth-y-reserva-modular https://github.com/Nadhe15/Atril-.git
cd Atril-

# 2. Instalar dependencias del proyecto
npm install

# 3. Ya la base de datos se encuentra conectada

# 4. Verificación estricta de tipos TypeScript (Resultado actual real: 0 errores - PASS)
npx tsc --noEmit

# 5. Ejecutar la aplicación en entorno local (Expo Web / Servidor de Desarrollo)
npx expo start --web

\end{lstlisting}
\end{tcolorbox}

\vspace{0.08cm}
\begin{tcolorbox}[
  colback=atrilWarnTint,
  colframe=atrilWarnText,
  boxrule=0.7pt,
  arc=3pt,
  left=7pt,
  right=7pt,
  top=4pt,
  bottom=4pt
]
  {\scriptsize\color{atrilWarnText}
  \textbf{Nota de Veracidad Técnica (Auditoría QA):} El chequeo estático de tipos \texttt{npx tsc --noEmit} finaliza con código \texttt{0} (sin errores en los 32 módulos TypeScript). La suite de pruebas unitarias/integración automatizadas (\texttt{\_\_tests\_\_}) y el porcentaje de cobertura con Jest/Vitest se dejan marcados como \textbf{pendientes por implementar} (\texttt{[PENDIENTE\_TESTS\_Y\_COBERTURA]}).}
\end{tcolorbox}

% --- Sección 4: Protocolo de Inyección de Defectos Deliberados ---
\secbadge{4}{Protocolo de Inyección de Defectos Deliberados (Rúbrica --- Tabla 5)}
{\footnotesize\color{atrilTextBody}
Para completar la \textbf{Tabla 5 (Pruebas con Defectos Provocados)}, se especifican dos defectos reales y reproducibles en los servicios y utilidades del repositorio:}

\vspace{0.10cm}

% Defecto A
\begin{tcolorbox}[
  colback=atrilDangerTint,
  colframe=atrilDanger,
  boxrule=0.8pt,
  arc=4pt,
  left=8pt,
  right=8pt,
  top=5pt,
  bottom=5pt
]
  {\footnotesize\bfseries\color{atrilDangerText} Defecto A (Nivel Unitario --- Lógica de Dominio y Validación en Servicios / Utilidades):}\\[2pt]
  {\footnotesize\color{atrilTextBody}
  \textbf{Archivo y función exacta:} \texttt{src/utils/escrowCalculator.ts} en \texttt{calcularHonorarioClase(tarifaBruta: number)} (líneas 22--37) \emph{[o en \texttt{src/services/authService.ts}, función \texttt{registrarUsuarioUnificado()}, línea 185]}.\\[1pt]
  \textbf{Acción provocada:} En \texttt{src/utils/escrowCalculator.ts}, comentar la validación de las líneas 23--25 (\texttt{if (!Number.isFinite(tarifaBruta) || tarifaBruta <= 0) throw new Error(...)}) o alterar la línea 30 cambiando \texttt{bruto - comisionPlataforma} por \texttt{bruto + comisionPlataforma}. \emph{(Alternativamente en \texttt{src/services/authService.ts}, línea 185: cambiar \texttt{const estadoInicial: EstadoVerificacion = 'pendiente'} a \texttt{'verificado'})}.\\[1pt]
  \textbf{Comando de prueba:} Ejecutar \texttt{npx tsc --noEmit} y evaluar \texttt{calcularHonorarioClase(0)} / \texttt{calcularHonorarioClase(25)} (o registrar un postulante en \texttt{AuthLandingScreen} con \texttt{npx expo start --web}).\\[1pt]
  \textbf{Resultado de fallo esperado:} La función admite tarifas nulas o negativas ($\le 0$) sin lanzar \texttt{Error("La tarifa bruta de la clase debe ser mayor a \$0.00 USD.")} y rompe la invariante contable ($\texttt{comisionPlataforma} + \texttt{honorarioNeto} \neq \texttt{bruto}$); o en \texttt{authService.ts} omite el bloqueo de \texttt{PendingVerificationScreen} otorgando acceso verificado sin auditoría de FundaOboe.}
\end{tcolorbox}

\vspace{0.10cm}

% Defecto B
\begin{tcolorbox}[
  colback=atrilSurfaceSub,
  colframe=atrilPrimary,
  boxrule=0.8pt,
  arc=4pt,
  left=8pt,
  right=8pt,
  top=5pt,
  bottom=5pt
]
  {\footnotesize\bfseries\color{atrilPrimary} Defecto B (Nivel de Integración --- Flujo de Verificación Admin y Restricción de Chat Interno):}\\[2pt]
  {\footnotesize\color{atrilTextBody}
  \textbf{Archivo y función exacta:} \texttt{src/services/chatService.ts} en \texttt{enviarMensaje(conversacionId, remitenteId, contenido)} (líneas 228--236) y \texttt{src/services/adminService.ts} en \texttt{listarUsuariosVerificacion(filtroEstado)} (líneas 166--220).\\[1pt]
  \textbf{Acción provocada:} (1) En \texttt{src/services/chatService.ts} (línea 234), remover la condición \texttt{|| !contenido.trim()} de la guardia de entrada; o (2) en \texttt{src/services/adminService.ts} (líneas 191 y 219), comentar el filtro \texttt{.filter((u) => u.estado\_verificacion === filtroEstado)} devolviendo \texttt{todosCache} sin filtrar.\\[1pt]
  \textbf{Comando de prueba:} Ejecutar \texttt{npx expo start --web} (o prueba de integración de servicios), invocar \texttt{enviarMensaje(convId, userId, "   ")} en \texttt{ChatRoomScreen} o seleccionar el filtro \emph{Pendientes} en \texttt{AdminDashboardScreen}.\\[1pt]
  \textbf{Resultado de fallo esperado:} El servicio de mensajería persiste mensajes en blanco en \texttt{public.mensajes} y \texttt{MENSAJES\_POR\_CONVERSACION\_CACHE} en lugar de lanzar \texttt{Error("La conversación, el remitente y el contenido del mensaje son obligatorios.")}, mientras que la consola Superadmin infringe el filtrado de estados exponiendo cuentas ya verificadas como pendientes.}
\end{tcolorbox}

% --- Sección 5: Cuadro de Conformidad y Firma de Coevaluación ---
\secbadge{5}{Cuadro de Conformidad y Firma de Coevaluación}

\begin{tcolorbox}[
  colback=atrilBg,
  colframe=atrilBorder,
  boxrule=0.9pt,
  arc=4pt,
  left=10pt,
  right=10pt,
  top=8pt,
  bottom=8pt
]
  \noindent
  \begin{tabularx}{\textwidth}{@{}p{0.44\textwidth}@{\hspace{12pt}}Y@{}}
    \begin{minipage}[t]{\linewidth}\footnotesize\color{atrilNavy}
      \vspace{2pt}
      \textbf{Equipo Evaluador:} \underline{\hspace{4.2cm}}\\[10pt]
      \textbf{Representante:} \underline{\hspace{4.6cm}}\\[10pt]
      \textbf{Fecha de Verificación:} \underline{\hspace{0.7cm}} / \underline{\hspace{0.7cm}} / 2026\\[10pt]
      \textbf{Firma de Conformidad:} \underline{\hspace{3.6cm}}
    \end{minipage}
    &
    \begin{minipage}[t]{\linewidth}\footnotesize\color{atrilTextBody}
      \begin{itemize}[leftmargin=0.5cm, itemsep=5pt, topsep=0pt, parsep=0pt]
        \item[\textcolor{atrilPrimary}{\large$\square$}] \textbf{Ejecución / Despliegue verificado correctamente} (Expo Web / Local y conexión Supabase BaaS).
        \item[\textcolor{atrilPrimary}{\large$\square$}] \textbf{Flujos de Roles (Admin, Profesor, Estudiante, Pendiente) y Chat verificados} en interfaz.
        \item[\textcolor{atrilPrimary}{\large$\square$}] \textbf{Verificación de tipos (\texttt{npx tsc --noEmit}) y tests automáticos ejecutados} (según Sección 3).
        \item[\textcolor{atrilPrimary}{\large$\square$}] \textbf{Defectos deliberados comprobados (Tabla 5)} (Defecto A Unitario y Defecto B Integración).
      \end{itemize}
    \end{minipage}
  \end{tabularx}
\end{tcolorbox}

\end{document}
